\sscp{Cuscus \it{(Phalangeridae)}}{C-03-01}
\citet{Schapper2011} reconstructs **wagal ‘cuscus’ which we transcribe
**waŋgal according to our interpretation of Proto-\ili{Aru} phonology (\srf{06-01}).
\citet{Nivens2017} also reconstructs Proto-South \ili{Aru} **yikuran ‘kind of cuscus’.
Reflexes of **mantəR in \ili{Aru} share with \ili{Fordata},
\ili{Kei}, and \ili{Banda} a local irregular development of penultimate *a
> **ə; **mantəR > **məntəR (\srf{14-04-03-02}).
\mbox{}\\

\noindent{\wg\st{2pt}\begin{tabularx}{\linewidth}
{lllll>{\raggedright\arraybackslash}X}
\lsptoprule
%Language	&	ISO	&	Term	&	Gloss	&	Etymon	&	Source	\\	\midrule
%\endfirsthead
Language	&	ISO	&	Term	&	Gloss	&	Etymon	&	Source	\\	\midrule
%\endhead
\ili{Ujir}	&	ugj	&	meday	&	cuscus	&	**məntəR	&	\cite{Schapper2011}	\\
West \ili{Kola}	&	kvv	&	medah;	&	cuscus,	&	**məntəR;	&	\cite{Schapper2011}	\\	\hhline{~}
	&		&	medah waŋal	&	male cuscus	&	**waŋgal	&	\hp{\,}	\\
East \ili{Kola}\footnotemark	&	kvv	&	modah, moduh; 	&	cuscus;	&	**məntəR;	&	\cite{Schapper2011};	\\	\hhline{~}
	&		&	moduh wagal;	&	male cuscus;	&	**waŋgal;	&	\cite{Nivens2017};	\\	\hhline{~}
	&		&	dilɸʷol	&	k.o. cuscus	&	\it{misc.}	&	\cite{Nivens2017}	\\
\ili{Lorang}	&	lrn	&	gwaŋal	&		&	**waŋgal	&	\cite{Schapper2011}	\\
\ili{Dobel}	&	kvo	&	kwaŋal;	&	cuscus;	&	**waŋgal;	&	\cite{Schapper2011}	\\	\hhline{~}
	&		&	iʔuran	&	k.o. cuscus	&	**yikuran	&	\cite{Nivens2017}	\\
\ili{Koba}	&	kpd	&	kwaŋal	&		&	**waŋgal	&	\cite{Schapper2011}	\\
\ili{Batuley}	&	bay	&	gwaŋal	&		&	**waŋgal	&	\cite{Schapper2011}	\\
\ili{Manombai}	&	woo	&	gaŋal	&		&	**waŋgal	&	\cite{Schapper2011}	\\
NW. \ili{Tarangan}	&	txn	&	ʤikuran	&	k.o. cuscus	&	**yikuran	&	\cite{Nivens2017}	\\
SW. \ili{Tarangan}	&	txn	&	medar;	&	cuscus;	&	**məntəR;	&	\cite{Schapper2011};	\\	\hhline{~}
	&		&	ʤikôran	&	k.o. cuscus	&	**yikuran	&	\cite{Nivens2017}	\\
\lspbottomrule
\end{tabularx}}
\footnotetext{\citet[264]{Schapper2011} gives \ili{Kola} \it{modah} ‘cuscus’.
		\citet{Nivens2017} has East \ili{Kola} \it{moduh wagal} ‘male cuscus’,
		\it{moduh bon} ‘female cuscus'
		and \it{moduh} ‘cuscus (generic)’.
		(All under the headword *wagal.) East \ili{Kola} \it{dilɸʷol} ‘k.o.
		cuscus’ occurs under the headword (*(sy)iku(Rr)an) for which
		\citet{Nivens2017} reconstructs Proto-South \ili{Aru} *yikuran.}

%\noindent{\wg\st{2pt}\begin{tabularx}{\linewidth}
%{lllll>{\raggedright\arraybackslash}X}
%\lsptoprule
%%Language	&	ISO	&	Term	&	Gloss	&	Etymon	&	Source	\\	\midrule
%%\endfirsthead
%Language	&	ISO	&	Term	&	Gloss	&	Etymon	&	Source	\\	\midrule
%%\endhead
%\ili{Ujir}	&	ugj	&	meday	&	cuscus	&	**məntəR	&	\cite{Schapper2011}	\\
%West \ili{Kola}	&	kvv	&	medah;	&	cuscus;	&	**məntəR;	&	\cite[264, 266]{Schapper2011}	\\	\hhline{~}
	%&		&	medah waŋal	&	male cuscus	&	**waŋgal	&	\hp{\,}	\\
%East \ili{Kola}	&	kvv	&	modah, moduh; 	&	cuscus;	&	**məntəR;	&	\cite{Schapper2011};	\\	\hhline{~}
	%&		&	moduh wagal;	&	male cuscus;	&	**waŋgal;	&	\cite{Nivens2017}	\\	\hhline{~}
	%&		&	dilɸʷol	&	k.o. cuscus	&	\it{misc.}	&	\hp{\,}	\\
%\ili{Lorang}	&	lrn	&	gwaŋal	&		&	**waŋgal	&	\cite{Schapper2011}	\\
%\ili{Dobel}	&	kvo	&	kwaŋal;	&	cuscus;	&	**waŋgal;	&	\cite{Schapper2011}	\\	\hhline{~}
	%&		&	iʔuran	&	k.o. cuscus	&	**yikuran	&	\cite{Nivens2017}	\\
%\ili{Koba}	&	kpd	&	kwaŋal	&		&	**waŋgal	&	\cite{Schapper2011}	\\
%\ili{Batuley}	&	bay	&	gwaŋal	&		&	**waŋgal	&	\cite{Schapper2011}	\\
%\ili{Manombai}	&	woo	&	gaŋal	&		&	**waŋgal	&	\cite{Schapper2011}	\\
%NW. \ili{Tarangan}	&	txn	&	ʤikuran	&	k.o. cuscus	&	**yikuran	&	\cite{Nivens2017}	\\
%SW. \ili{Tarangan}	&	txn	&	medar;	&	cuscus;	&	**məntəR;	&	\cite[263]{Schapper2011};	\\	\hhline{~}
	%&		&	ʤikôran	&	k.o. cuscus	&	**yikuran	&	\cite{Nivens2017}	\\
%\lspbottomrule
%\end{tabularx}}

%\noindent{\wg\st{2pt}\begin{xltabular}{\linewidth}
%{lllll>{\raggedright\arraybackslash}X}
%\lsptoprule
%%Language	&	ISO	&	Term	&	Gloss	&	Etymon	&	Source	\\	\midrule
%%\endfirsthead
%Language	&	ISO	&	Term	&	Gloss	&	Etymon	&	Source	\\	\midrule
%\endhead
%
%\lspbottomrule
%\end{xltabular}}